% !PRINTED 227
\weberpage{227}
%<<CONT
Sind $m, n$ irgend zwei teilerfremde Zahlen, letztere ungerade, bedeutet ferner $\mu$ die Anzahl derjenigen unter den Zahlen

\[
m,\ 2m,\ 3m \cdots \frac{n-1}{2}\,m,
\]

deren absolut kleinste Reste $(\mathrm{mod}\ n)$ negativ sind, so ist

\weq{8}{(-1)^\mu = \left(\frac{m}{n}\right),}

worin $\left(\dfrac{m}{n}\right)$ das \gesperrt{Legendre-Jacobische} Symbol aus der Theorie der quadratischen Reste ist\textsuperscript{1)}.

Dieser Satz führt nun unmittelbar zur Bestimmung des Vorzeichens in (7). Denn wenn in (5) die Funktion $\Phi(x)$ ungerade ist, so ändern beim Übergang von $\Omega$ zu $m\,\Omega$ genau $\mu$ Faktoren in (5) ihr Vorzeichen und wir schließen

\weq{9}{\Pi(m\,\Omega) = \left(\frac{m}{n}\right) \Pi(\Omega).}

Ist $n$ keine Quadratzahl, so kann man $m$ immer so annehmen, daß $\left(\dfrac{m}{n}\right) = -1$ ist. Ist nämlich $n = p^\lambda n'$, $\lambda$ ungerade, $n'$ nicht durch $p$ teilbar, $\beta$ ein quadratischer Nichtrest von $p$, so braucht man $m$ nur aus den Kongruenzen

\[
m \equiv \beta \pmod{p^\lambda}, \quad m \equiv 1 \pmod{n'}
\]

zu bestimmen, um eine solche Zahl $m$ zu finden. Und dann ist $\Pi(m\,\Omega) = -\Pi(\Omega)$ und nicht $\Pi(\Omega)$, sondern erst $\Pi(\Omega)^2$ Wurzel einer Transformationsgleichung. Ist aber $n$ eine Quadratzahl, so ist $\Pi(m\,\Omega) = \Pi(\Omega)$ für jedes $m$ und folglich $\Pi(\Omega)$ selbst Wurzel einer Transformationsgleichung. Es folgt also der Satz:

\begin{quote}
\gesperrt{Bei ungerader Funktion $\Phi$ ist $\Pi(\Omega)^2$, und nur wenn $n$ eine Quadratzahl ist, $\Pi(\Omega)$ selbst, Wurzel einer Transformationsgleichung\textsuperscript{2)}.}
\end{quote}

\textsuperscript{1)}~Vgl.\ über diesen Satz: \gesperrt{Schering} und \gesperrt{Kronecker} im Monatsbericht der Berliner Akademie vom 22.\,Juni 1876. \gesperrt{Schering}, Acta mathematica I.

\textsuperscript{2)}~Diese Vereinfachung der Multiplikatorgleichung in dem Fall, wo $n$ ein Quadrat ist, hat \gesperrt{Joubert} entdeckt, aber auf einem von dem unserigen ganz verschiedenen Wege nachgewiesen. \glqq Sur les équations, qui se rencontrent dans la théorie de la transformation des fonctions elliptiques.\grqq{} Paris 1876.
